\documentclass[10pt, openany]{book}
%
%  Packages to use
%
\usepackage{fancyhdr}
\usepackage{fancyvrb}
\usepackage{fancybox}
%
\usepackage{lastpage}
\usepackage{imakeidx}
%
\usepackage{amsmath}
\usepackage{amsfonts}
%
\usepackage{geometry}
\geometry{letterpaper}
%
\usepackage{url}
\usepackage{gensymb}
\usepackage{multicol}
\usepackage{xcolor}
%
\usepackage[pdf]{pstricks}
\usepackage{graphicx}
\DeclareGraphicsExtensions{.pdf}
\DeclareGraphicsRule{.pdf}{pdf}{.pdf}{}
%
% Rules to allow import of graphics files in EPS format
%
\usepackage{graphicx}
\DeclareGraphicsExtensions{.eps}
\DeclareGraphicsRule{.eps}{eps}{.eps}{}
%
%  Include the listings package
%
\usepackage{listings}
%
%  Define Tiny Lisp based on Common Lisp
%
\lstdefinelanguage[Tiny]{Lisp}[]{Lisp}{morekeywords=[13]{atomp, bit-vector-p, car, cdr, char-downcase, char-code, char-upcase, compiled-function-p, dowhile, dump, exit, fresh-line, if, code-char, lambda, msg, nullp, parse-integer, peek8, peek16, peek32, poke8, poke16, poke32, progn, quote, read-line, reset, setq, simple-bit-vector-p, simple-string-p, simple-vector-p, string-downcase, string-upcase, memb, meml, memw, sim-step, go, reg-val, num-reg, halted}}
\lstset{language=Ada}
%
% Macro definitions
%
\newcommand{\operation}[1]{\textbf{\texttt{#1}}}
\newcommand{\package}[1]{\texttt{#1}}
\newcommand{\function}[1]{\texttt{#1}}
\newcommand{\constant}[1]{\emph{\texttt{#1}}}
\newcommand{\keyword}[1]{\texttt{#1}}
\newcommand{\datatype}[1]{\texttt{#1}}
\newcommand{\comment}[1]{{\color{red}{#1}}}
%
\newcommand{\docname}{Users's Manual for CPU Simulators \\  CLI and Lisp Programming}

%
% Front Matter
%
\title{\docname}
\author{Brent Seidel \\ Phoenix, AZ}
\date{ \today }
%========================================================
%%% BEGIN DOCUMENT
\begin{document}
%
%  Header's and Footers
%
\fancypagestyle{plain}{
  \fancyhead[L]{}%
  \fancyhead[R]{}%
  \fancyfoot[C]{Page \thepage\ of \pageref{LastPage}}%
  \fancyfoot[L]{CPU Simulation CLI and Lisp}
  \renewcommand{\headrulewidth}{0pt}%
  \renewcommand{\footrulewidth}{0.4pt}%
}
\fancypagestyle{myfancy}{
  \fancyhead[L]{\docname}%
  \fancyhead[R]{\leftmark}
  \fancyfoot[C]{Page \thepage\ of \pageref{LastPage}}%
  \fancyfoot[L]{CPU Simulation CLI and Lisp}
  \renewcommand{\headrulewidth}{0.4pt}%
  \renewcommand{\footrulewidth}{0.4pt}%
}
\pagestyle{myfancy}
%
% Produce the front matter
%
\maketitle
\begin{center}
This document is \copyright 2026 Brent Seidel.  All rights reserved.

\paragraph{}Note that this is a draft version and not the final version for publication.
\end{center}
\tableofcontents

%========================================================
\chapter{Introduction}
This project is for a CLI interface to the CPU simulator project.
In addition to using the simulations this way, it can also serve as an example for embedding a simulator in another application.
For more information on the simulations themselves and their internals, see the simulations document.

\section{Disclaimer}
I have tried to make the documentation and software as accurate as possible.
However, it is inevitable that errors exist.
Should you find an error, feel free to write an issue on the GitHub project and I will look at it.
For software errors especially, please provide information as to what change is needed.
If you would like some particular hardware added, feel free to request that as well.
Note that this is a hobby project maintained by one person who has other demands on their time.

\section{License}
This project is licensed using the GNU General Public License V3.0.
Should you wish other licensing terms, contact the author.

THE SOFTWARE IS PROVIDED "AS IS", WITHOUT WARRANTY OF ANY KIND, EXPRESS OR IMPLIED, INCLUDING BUT NOT LIMITED TO THE WARRANTIES OF MERCHANTABILITY, FITNESS FOR A PARTICULAR PURPOSE AND NONINFRINGEMENT. IN NO EVENT SHALL THE AUTHORS BE LIABLE FOR ANY CLAIM, DAMAGES OR OTHER LIABILITY, WHETHER IN AN ACTION OF CONTRACT, TORT OR OTHERWISE, ARISING FROM, OUT OF OR IN CONNECTION WITH THE SOFTWARE OR THE USE OR OTHER DEALINGS IN THE SOFTWARE.

%========================================================
\chapter{How To Obtain}
This package is currently available on GitHub at \url{https://github.com/BrentSeidel/Sim-CPU}

\section{Dependencies}
\subsection{Ada Libraries}
The following Ada packages are used:
\begin{itemize}
  \item \package{Ada.Calendar}
  \item \package{Ada.Characters.Latin\_1}
  \item \package{Ada.Containers.Indefinite\_Ordered\_Maps}
  \item \package{Ada.Direct\_IO}
  \item \package{Ada.Exceptions}
  \item \package{Ada.Integer\_Text\_IO}
  \item \package{Ada.Streams}
  \item \package{Ada.Strings.Fixed}
  \item \package{Ada.Strings.Maps.Constants}
  \item \package{Ada.Strings.Unbounded}
  \item \package{Ada.Text\_IO}
  \item \package{Ada.Text\_IO.Unbounded\_IO}
  \item \package{Ada.Unchecked\_Conversion}
\end{itemize}


\subsection{Other Libraries}
This library depends on the root package \package{BBS} available at \url{https://github.com/BrentSeidel/BBS-Ada} or using alire via ``\keyword{alr get bbs}'' and the \package{BBS.Lisp} package available at \url{https://github.com/BrentSeidel/Ada-Lisp} or using alire via ``\keyword{alr get bbs\_lisp}''.

%========================================================
\chapter{Usage Instructions}


\section{Using Alire}
Alire automatically handles dependencies.
To use this in your project, just issue the command ``\keyword{alr with bbs\_simcpu}'' in your project directory.
To build the standalone CLI program, first obtain the cli using ``\keyword{alr get simcpucli}''.
Change to the appropriate directory and use ``\keyword{alr build}'' and ``\keyword{alr run}''.
To use the load CP/M program, use ``\keyword{alr get loadcpm}''.
Change to the appropriate directory and use ``\keyword{alr build}'' and ``\keyword{alr run}''.

\section{Using \keyword{gprbuild}}
This is a library of routines intended to be used by some program.
To use these in your program, edit your *\keyword{.gpr} file to include a line to \keyword{with} the path to \keyword{bbs\_simcpu\_noalr.gpr}.
Then in your Ada code \keyword{with} in the package(s) you need and use the routines.

This is also available as a standalone program with a command line interface (CLI).
This can be used for development or debugging.
To build the CLI, just build the \keyword{simcli\_noalr.gpr} project file, after making sure that the dependencies are in the expected places.

There is another program \keyword{loadcpm} available to create a bootable CP/M floppy disk image.
It is build using the \keyword{loadcpm\_noalr.gpr} project.
It takes an Intel hex file image and writes it to a floppy disk image.

%========================================================
\chapter{General}
\section{Available Simulators}
Several simulators are available for use.  Each simulator may also have variation.  So, one simulator may provide variations for different processors in a family of processors.


%--------------------------------------------------------------------------------------------------
\subsection{8080 Family}
The 8080 simulator has variants defined for the 8080, 8085, and Z-80 processors.
These are 8 bit processors with an 8 bit data bus and a 16 bit address bus.
In addition to a memory bus, these processors also include an I/O bus with 8 bit I/O port addressing.
See \cite{i8080} for more information about the 8080/8085.
See \cite{z80} and \cite{z80undoc} for more information about the Z-80.

Currently, the 8080 family does not have memory mapped I/O.  This may be added at some time in the future.

When loading data using the \function{load} routine, the specified file is assumed to be in Intel Hex format.

%--------------------------------------------------------------------------------------------------
\subsection{68000 Family}
The 68000 simulator has variants defined for the 68000, 68008, 68010, and CPU32.
Only the 68000 and 68008 are currently implemented.
The 68010 and CPU32 are for future development.
Internally, these are 32 bit processors with 32 bit data and 32 bit address busses.
These processors do not have a separate I/O address space.
All I/O is mapped as part of the main memory.
The external address and data bus sizes depend on the variant selected.
See \cite{m68000a} and \cite{m68000b} for more information about the 68000 family.

The interrupt code is interpreted with the low order bits (7-0) representing the vector number and the next 8 bits (15-8) representing the priority.
Interrupts with vectors that match internally defined exception vectors are ignored.
Thus only interrupt vectors in the range 25-31 and 64-255 are processed.

When loading data using the \function{load} routine, the specified file is assumed to be in Motorola S-Record format.

%--------------------------------------------------------------------------------------------------
\subsection{6502 Family}
The 6502 simulator currently has only a working variants defined for the 6502 processor.
This is an 8 bit processors with an 8 bit data bus and a 16 bit address bus.
This processor does not have a separate I/O address space.
All I/O is mapped as part of the main memory.
Three interrupts are define for normal interrupts, non-maskable interrupts, and reset.
See \cite{msc6502} for more information about the 6502.

When loading data using the \function{load} routine, the specified file is assumed to be in Intel Hex format.

%--------------------------------------------------------------------------------------------------
\subsection{PDP-11 Family}
Most of the 16 bit addressing members are supported, except the LSI-11, PDP-11/03.
Future models with memory management are planned at some point.
This is a line of 16 bit minicomputers that was developed in the 1970s and was popular through the 1990s.
Models with no memory management had 16 address bit.
With memory management, depending on the model, addresses of 18 or 22 bits was supported.
These processors do not have a separate I/O address space.
All I/O is mapped as part of the main memory.
By convention, all I/O is located in the upper 8kBytes of memory.
See \cite{dec80} and \cite{dec84} for more information about the PDP-11.

A wealth of information about the PDP-11 (and other DEC hardware) is on bitsavers at \url{https://www.bitsavers.org/pdf/dec/}.
Software can also be found on bitsavers at \url{https://www.bitsavers.org/bits/DEC/}.
Various diagnostic software can be found using the PDP-11 Diagnostics - Database at \url{https://www.retrocmp.com/tools/pdp-11-diagnostic-database/202-pdp-11-diagnostics-database}.

When loading data using the \function{load} routine, the specified file is assumed to be in DEC Absolute Loader format.

%========================================================
\section{I/O Devices}
\label{chap:IODev}
A number of simulated I/O devices are provided.
Some are designed to mimic real I/O devices while others are just for the simulator.

Each device has a name that can be used to refer to it.
The name consists of an alphabetic device code followed by a decimal unit number.
If no unit number is specified, 0 is assumed.  A unit number of zero is special.
It is used to refer to a generic device, such as when attaching a device - the unit number is assigned during the attachment.
The currently defined device codes are:
\begin{description}
  \item[BM792] -- A boot ROM for booting a PDP-11 from a RK11 controller.
  \item[CLK] -- A clock device to generate a periodic interrupt.
  \item[FD] -- Main disk controller with a track and sector interface.  It started out as a floppy disk controller, but geometries can be defined beyond that.
  \item[MUX] -- An 8 channel terminal multiplexer using a telnet interface.
  \item[PRN] -- A simple printer output device.
  \item[PTP] -- A paper tape interface.  It can read and write simulated tapes.
  \item[TEL] -- A single channel terminal device using a telnet interface.
\end{description}
The following devices simulate actual hardware for PDP-11 computers.
\begin{description}
  \item[DL11] --  A single line serial interface, see \cite{DL11}.
  \item[DZ11] -- An 8 line serial interface multiplexer, see \cite{DZ11}.
  \item[KT11] -- A memory management unit for the PDP-11, see \cite{KT11}.
  \item[KW11] -- A clock for the PDP-11 series of computer, see \cite{KW11L}.
  \item[PC11] -- A combination paper tape reader/punch, see \cite{PC11}.
  \item[RK11] -- A disk controller that can have up to 8 attached RK05 simulated drives, see \cite{RK11}.
  \item[RK611] -- A disk controller than can have up to 8 attached RK07 simulated drives (RK06 drives are not supported), see \cite{RK611}.
  \item[TM11] -- A magnetic tape controller that can have up to 8 attached tape drives, see \cite{TM11}.
\end{description}

%========================================================
\chapter{Command Line Interface}
%--------------------------------------------------------------------------------------------------
\section{Commands}
A number of commands are provided to allow the user to run and interact with software running on the simulator.
The commands and their exact operation may change depending on feedback from usage.
A dash (`-') in the name is used to indicate the abbreviation point for the command.
\begin{description}
  \item[;] -- A comment.  Everything following on the line is ignored, except if the rest of the line is being used as a file name.
  \item[ATT-ACH $<$device$>$ $<$addr$>$ $<$bus$>$ [$<$dev-specific$>$]] --  Attaches a specific I/O device at the bus address.  Depending on the device, device specific parameters may be required.
  \item[BENC-HMARK] --  Similar to \verb|RUN|, except that it measures the number of instructions executed per second.
  \item[BREA-K $<$addr$>$] -- Sets a breakpoint at the specified address.
  \item[C-ONTINUE] -- Clear a simulation halted condition.
  \item[CPU $<$cpu-name$>$] -- Set the CPU to be simulated.  This can only be done once per session and must be done before any command that requires a CPU.
  \item[DEP-OSITE $<$addr$>$ $<$byte$>$] -- Deposit a byte into memory.
  \item[DISK]  --  Commands related to floppy disks and image files (see below for more details).
  \item[D-UMP $<$addr$>$] -- Dump a block of memory starting at the specified address..
  \item[EXIT] -- Exit the simulation (synonym for \verb|QUIT|).
  \item[GO $<$addr$>$] -- Sets the execution address.
  \item[INTERRUPT]-- Can be abbreviated to \verb|INT|.  Used to enable or disable interrupt processing by the simulated CPU.
  \item[LISP [$<$filename$>$]] -- Starts the Tiny-Lisp environment.  If a filename is specified, Lisp commands come from that file.
  \item[LIST [$<$devname$>$]] -- Lists all the I/O devices attached to the simulator.  If $<$devname$>$ is supplied, it provides information for that specific device.
  \item[LOAD $<$filename$>$] -- Loads memory with the specified file (typically Intel Hex or S-Record format).
  \item[PR-INT] -- Used for configuring printers.
  \item[TAPE] -- Commands relating to printer and attached file (see below for more details).
  \item[QU-IT] -- Exit the simulation (synonym for \verb|EXIT|).
  \item[REG-ISTER] -- Display the registers and their contents.
  \item[RES-ET] -- Calls the simulator's \verb|init| routine.
  \item[R-UN] -- Runs the simulator.  While running, it can be interrupted by a control-E character.
  \item[SET $<$subcommand and options$>$] -- Sets certain simulator options.
  \item[SH-OW $<$subcommand and options$>$] --  Shows certain simulator options.
  \item[STEP] -- Execute a single instruction, if the simulation is not halted.
  \item[TAPE] -- Commands relating to paper (or cassette) tape and attached files (see below for more details).
  \item[TRA-CE $<$value$>$] -- Sets the trace value.  The value is interpreted by the simulator and typically causes certain information to be printed while processing.
  \item[UNBR-EAK $<$addr$>$] -- Clears a breakpoint at the specified address.  For some simulators, the address is ignored and may not need to be specified.  In this case, all or the only breakpoint is cleared.
\end{description}

\subsection{Attach Command}
\label{cli:Attach}
The attach command is used to attach a device to the current CPU simulator.
The device is specified as a generic device name.
The address is given in decimal.
The bus is:
\begin{description}
  \item[IO] -- Use the I/O bus (only available on some CPUs).
  \item[MEM] -- Use memory mapped I/O (only available on some CPUs).
\end{description}
Some of the devices take additional parameters.
\begin{description}
  \item[BM792] --  Takes no additional options.
  \item[CLK] -- Takes one additional decimal value for the exception code to be presented to the CPU when an interrupt occurs.
  \item[DL11] -- Takes two additional decimal parameters.  The first is the TCP/IP port for the first terminal (the additional terminals are assigned sequential port numbers).  The second is a 32 bit integer combination of two exception codes to be presented to the CPU when an interrupt occurs.  The lower order 16 bits are the receive vector and the upper 16 bits are the transmit vector.
  \item[DZ11] -- Takes two additional decimal parameters.  The first is the TCP/IP port for the first terminal (the additional terminals are assigned sequential port numbers).  The second is a 32 bit integer combination of two exception codes to be presented to the CPU when an interrupt occurs.  The lower order 16 bits are the receive vector and the upper 16 bits are the transmit vector.
  \item[FD] -- Takes one additional decimal parameter for the number of attached drives in the range 0-15, though 0 drives is rather pointless.
  \item[KT11] -- Takes one additional decimal value for the exception code to be presented to the CPU when an interrupt occurs.
  \item[KW11] -- Takes one additional decimal value for the exception code to be presented to the CPU when an interrupt occurs.
  \item[MUX] -- Takes two additional decimal parameters.  The first is the TCP/IP port for the first terminal (the additional terminals are assigned sequential port numbers).  The second is the exception code to be presented to the CPU when an interrupt occurs.
  \item[PC11] -- Takes one additional parameter.  It is a 32 bit integer combination of two exception codes to be presented to the CPU when an interrupt occurs.  The lower order 16 bits are the receive vector and the upper 16 bits are the transmit vector.
  \item[PRN] -- Takes no additional parameters.
  \item[PTP] -- Takes no additional parameters.
  \item[RK11] -- Takes one additional decimal value for the exception code to be presented to the CPU when an interrupt occurs.
  \item[RK611] -- Takes one additional decimal value for the exception code to be presented to the CPU when an interrupt occurs.
  \item[TEL] --  Takes two additional decimal parameters.  The first is the TCP/IP port for the first terminal (the additional terminals are assigned sequential port numbers).  The second is the exception code to be presented to the CPU when an interrupt occurs.
  \item[TM11] -- Takes one additional decimal value for the exception code to be presented to the CPU when an interrupt occurs.
\end{description}

\subsection{CPU Command}
This command is deprecated.  SET CPU is the preferred command.

\subsection{Disk Commands}
A number of subcommands are available for interfacing with simulated floppy disks and image files.
The DISK command takes the form:
\begin{itemize}
  \item DISK $<$controller$>$ $<$subcommand$>$ $<$parameters for subcommand$>$
\end{itemize}
The controller is the device name of an attached disk controller.  The subcommands are:
\begin{description}
  \item[CLOSE $<$drive number$>$] -- Closes the image file associated with the specified drive number.
  \item[OPEN $<$drive number$>$ $<$file name$>$] -- Attaches the specified file to the specified drive.
  \item[READONLY $<$drive number$>$] -- Sets the specified drive to read-only.
  \item[READWRITE $<$drive number$>$] -- Sets the specified drive to read-write.
\end{description}

\subsection{Interrupt Commands}
A number of subcommands are available for interfacing with interrupts.
\begin{description}
  \item[ON] -- Enables interrupt processing by the simulator (if supported).
  \item[OFF] -- Disables interrupt processing by the simulator.  This may be useful for single-stepping through some code without getting interrupted by the clock or other devices.
  \item[SEND $<$interrupt number$>$] --  Sends the specified interrupt number to the simulator (if supported).
\end{description}

\subsection{Show subcommands}
\label{cmd-show}
This allows viewing of some simulator options.
The current options are:
\begin{description}
  \item[OP-TIONS $<$option-name$>$] -- Shows the value of a CPU specific option.
  \item[TR-ACE] -- Shows the current trace settings
\end{description}

\subsection{Set Commands}
\label{cmd-set}
This allows setting of some simulator options.
The current options are:
\begin{description}
  \item[CPU] -- The same as the CPU command above.  The following CPUs are currently defined:
  \begin{description}
    \item[8080] -- The Intel 8080 processor.  An 8 bit processor that can address up to 64kBytes of memory and 256 I/O ports.  See \cite{i8080}.
    \item[8085] -- Made by Intel.  Very similar to the 8080 with mostly hardware enhancements, with 2 added instructions.  See \cite{i8080}.
    \item[Z80] -- Made by Zilog.  It is mostly upward compatible with the 8080 (one flag is different on some instructions).  The added instruction codes for the 8085 are completely different instructions on the Z80.  See \cite{z80} and \cite{z80undoc}.
    \item[68000] -- The Motorola 68000 processor.  This is internally a 32 bit processor with a 16 bit data bus, and 24 address lines.  It can address up to 16MBytes of memory.  See \cite{m68000a} and \cite{m68000b}.
    \item[68008] -- A stripped down version of the 68000.  The core is the same, but it has an 8 bit address bus and only 20 address lines.  See \cite{m68000a} and \cite{m68000b}.
    \item[6502] -- A simple 8 bit processor with 16 bits of addressing.  See \cite{msc6502}.
    \item[PDP-11/04] -- A version of the PDP-11 series of minicomputers.  See \cite{dec75} and \cite{dec80}.
    \item[PDP-11/05] -- A version of the PDP-11 series of minicomputers.  See \cite{PDP1105}, \cite{dec75} and \cite{dec80}.
    \item[PDP-11/10] -- The same as the PDP-11/05.
    \item[PDP-11/15] -- A version of the PDP-11 series of minicomputers.
    \item[PDP-11/20] -- The same as the PDP-11/15.
    \item[PDP-11/34] -- Under development.
    \item[PDP-11/35] -- Under development.  See \cite{dec75}.
    \item[PDP-11/40] -- The same as PDP-11/35.
  \end{description}
  It is likely that more PDP-11 models will be added.
  \item[OP-TIONS $<$option-name$>$ $<$option-value$>$] -- Sets a specified CPU option to the specified value.  The options and values are specific to a particular CPU.
  \item[PAUSE] -- Sets either the pause count or the pause character.  This controls how often to check for the pause/interrupt character while the simulator is running.  Since the check for the character takes much longer than simulating an instruction, increasing this value can speed the simulation.
\end{description}

\subsection{Tape Commands}
The paper tape interface consists of a reader called \verb|RDR| and a writer called \verb|PUN| (for punch).
Files can be opened or closed for either of these devices.
The magnetic tape interface may have multiple tape drives attached, depending on the controller.
The TAPE command takes the form:
\begin{itemize}
  \item TAPE $<$controller$>$ $<$subcommand$>$ $<$parameters for subcommand$>$
\end{itemize}
The controller is the device name of an attached tape controller.
The subcommands for paper tape are:
\begin{description}
  \item[CLOSE $<$PUN or RDR$>$]  -- Closes the file attached to the reader or punch.
  \item[OPEN $<$PUN or RDR$>$ $<$file name$>$] --  Attaches the specified file to the reader or punch.
\end{description}

The subcommands for magnetic tape are:
\begin{description}
  \item[CLOSE $<$drive number$>$] -- Closes the image file associated with the specified drive number.
  \item[OPEN $<$drive number$>$ $<$file name$>$] -- Attaches the specified file to the specified drive.
  \item[READONLY $<$drive number$>$] -- Sets the specified drive to read-only.
  \item[READWRITE $<$drive number$>$] -- Sets the specified drive to read-write.
\end{description}

\subsection{Printer Commands}
The printer device simply writes data sent to it to an output file.
If no file is open, the data is discarded.
The PRINT command takes the form:
\begin{itemize}
  \item PRINT $<$controller$>$ $<$subcommand$>$ $<$parameters for subcommand$>$
\end{itemize}
The controller is the device name of an attached printer controller.  The subcommands are:
\begin{description}
  \item[CLOSE]  -- Closes the file attached to the printer.
  \item[OPEN  $<$file name$>$] --  Attaches the specified file to the printer.
\end{description}

%--------------------------------------------------------------------------------------------------
\section{Lisp Programming}
The Lisp environment is based on Ada-Lisp, available at \url{https://github.com/BrentSeidel/Ada-Lisp}, with a few commands added to interface with a simulator.
Refer to the Ada-Lisp documentation for details about the core language.
This section just describes the additional commands.

\subsection{Attach}
\subsubsection{Inputs}
The following inputs are expected:
\begin{itemize}
  \item A \datatype{String} representing the device name.
  \item An \datatype{Integer} for the device address.
  \item A \datatype{String} representing the address bus used (``IO'' or ``MEM'').
  \item Additional items specific to the device may be required.
\end{itemize}
\subsubsection{Outputs}
None.
\subsubsection{Description}
Attaches an I/O device at the specified address and address bus.
Refer to the Attach command (section \ref{cli:Attach}) for more information.

\subsection{Disk-Close}
The following inputs are expected:
\begin{itemize}
  \item A \datatype{String} representing the disk controller name.
  \item An \datatype{Integer} for the drive number.
\end{itemize}
\subsubsection{Outputs}
None.
\subsubsection{Description}
Closes the file attached to the specified disk drive.

\subsection{Disk-Geom}
The following inputs are expected:
\begin{itemize}
  \item A \datatype{String} representing the disk controller name.
  \item An \datatype{Integer} for the drive number.
  \item A \datatype{String} representing the geometry (currently ``IBM'' for 8 inch diskettes or ``HD'' for a 5MB hard disk).  It is planned to extend this to allow a list specifying arbitrary geometry as well.  Note that some controllers may not allow setting of geometry or may otherwise restrict the values.
\end{itemize}
\subsubsection{Outputs}
None.
\subsubsection{Description}
Sets the disk geometry (tracks, sectors, heads) of the specified disk drive.

\subsection{Disk-Open}
The following inputs are expected:
\begin{itemize}
  \item A \datatype{String} representing the disk controller name.
  \item An \datatype{Integer} for the drive number.
  \item A \datatype{String} representing the file name to attach to the disk drive.
\end{itemize}
\subsubsection{Outputs}
None.
\subsubsection{Description}
Opens a file and attaches it to the specified disk drive.

\subsection{Disk-Protect}
The following inputs are expected:
\begin{itemize}
  \item A \datatype{String} representing the disk controller name.
  \item An \datatype{Integer} for the drive number.
  \item A \datatype{Boolean} with \constant{True} setting the drive to read-only.
\end{itemize}
\subsubsection{Outputs}
None.
\subsubsection{Description}
Opens a file and attaches it to the specified disk drive.

\subsection{Go}
\subsubsection{Inputs}
An \datatype{Integer} representing the starting address for program execution.
\subsubsection{Outputs}
None.
\subsubsection{Description}
Sets the starting address for program execution.
The next \verb|(sim-step)| command executes the instruction at this location.

\subsection{Halted}
\subsubsection{Inputs}
An optional \datatype{Boolean} that if set \constant{True}, clears the halted state of the simulator.
If set \constant{False}, does nothing.
\subsubsection{Outputs}
If no input is provided, it returns a \datatype{Boolean} representing the halted state, 
\subsubsection{Description}
This is used to clear to test the halted state of a simulator.

\subsection{Int-state}
\subsubsection{Inputs}
None.
\subsubsection{Outputs}
A simulator dependent \datatype{Integer} representing the state of interrupts.
\subsubsection{Description}
Returns a simulator dependent \datatype{Integer} representing the state of the interrupts.
Typically it is something like 0 for interrupts disabled and 1 for interrupts enabled.
For systems with priorities, it may be the processor priority.

\subsection{Last-out-addr}
\subsubsection{Inputs}
None.
\subsubsection{Outputs}
The address used by the last output instruction as an \datatype{Integer}.
\subsubsection{Description}
Returns the address used by the last output instruction.
This is only meaningful if the processor being simulated has a separate output bus rather than memory mapped I/O.

\subsection{Last-out-data}
\subsubsection{Inputs}
None.
\subsubsection{Outputs}
The data used by the last output instruction as an \datatype{Integer}.
\subsubsection{Description}
The data used by the last output instruction.
This is only meaningful if the processor being simulated has a separate output bus rather than memory mapped I/O.

\subsection{Memb}
\subsubsection{Inputs}
A memory address and an optional byte value, both as \datatype{Integer}s.
\subsubsection{Outputs}
If no byte value is provided, it returns the byte at the memory address as an \datatype{Integer}.
\subsubsection{Description}
Reads or writes a byte value in memory.

\subsection{Meml}
\subsubsection{Inputs}
A memory address and an optional long value, both as \datatype{Integer}s.
\subsubsection{Outputs}
If no long value is provided, it returns the long at the memory address as an \datatype{Integer}.
\subsubsection{Description}
Reads or writes a long value in memory.

\subsection{Memw}
\subsubsection{Inputs}
A memory address and an optional word value, both as \datatype{Integer}s.
\subsubsection{Outputs}
If no word value is provided, it returns the word at the memory address as an \datatype{Integer}.
\subsubsection{Description}
Reads or writes a word value in memory.

\subsection{Mem-limit}
\subsubsection{Inputs}
An optional \datatype{Integer} representing the value to set the memory limit to.
If this value is greater than (mem-max), it is ignored.
\subsubsection{Outputs}
If the input is not specified, returns the available memory size.
\subsubsection{Description}
Allows dynamically setting the maximum memory size to less than the configured memory size.
This is typically done for test purposes.

\subsection{Mem-max}
\subsubsection{Inputs}
None.
\subsubsection{Outputs}
An \datatype{Integer} representing the maximum memory size.
\subsubsection{Description}
Returns the maximum configured memory size for the simulated processor.

\subsection{Num-reg}
\subsubsection{Inputs}
None
\subsubsection{Outputs}
The number of registers defined by the simulator as an \datatype{Integer}.
\subsubsection{Description}
Returns the number of registers defined by the simulator.
This could be used for iterating through the registers and displaying their values.

\subsection{Option}
\subsubsection{Inputs}
Two \datatype{String} inputs, with the second being optional.
The first is the name of the option.
The optional second is the value to set the option to.
\subsubsection{Outputs}
If the second input is not specified, returns the value of the option as a \datatype{String}.
\subsubsection{Description}
Allows CPU specific options to be set and read.

\subsection{Override-in}
\subsubsection{Inputs}
Two \datatype{Integer}s representing the address of an I/O port and the data to provide.
\subsubsection{Outputs}
None
\subsubsection{Description}
Provides an address and data to override the next input instruction.
If the instruction reads from the address, the provided data is read by the simulator instead of reading from the actual I/O device.
Inputs from a different address work normally and any input instruction clears the override.
This is intended for testing purposes.

\subsection{Print-Close}
The following inputs are expected:
\begin{itemize}
  \item A \datatype{String} representing the printer controller name.
\end{itemize}
\subsubsection{Outputs}
None
\subsubsection{Description}
Closes the file attached to the specified printer.

\subsection{Print-Open}
The following inputs are expected:
\begin{itemize}
  \item A \datatype{String} representing the printer controller name.
  \item A \datatype{String} representing the file name to attach to the printer.
\end{itemize}
\subsubsection{Outputs}
None
\subsubsection{Description}
Opens a file and attaches it to the specified printer.

\subsection{Reg-val}
\subsubsection{Inputs}
A register number as an \datatype{Integer}.
\subsubsection{Outputs}
The value of the specified register as an \datatype{Integer}, or an error if the register number is not in range.
\subsubsection{Description}
Returns the value of the specified register.

\subsection{Send-int}
\subsubsection{Inputs}
An \datatype{Integer} interrupt code.
\subsubsection{Outputs}
None.
\subsubsection{Description}
Sends the specified interrupt code to the simulator for processing.
The processing done depends on the specific simulator.
It may be ignored, be a vector number, be an instruction, or something else.

\subsection{Set-pause-char}
\subsubsection{Inputs}
An \datatype{Integer} in the range 0-255 as the ASCII value of the pause character.
Note that if you're on a system that doesn't use ASCII, then this is whatever the numerical value for the character is.
\subsubsection{Outputs}
The current pause character.
\subsubsection{Description}
Sets the pause character.
Pressing this character while a simulation is running will cause the simulation to return to the command line.
If no character is specified, it just returns the current pause character.

\subsection{Set-pause-count}
\subsubsection{Inputs}
A positive \datatype{Integer}.
\subsubsection{Outputs}
The current value of the pause count.
\subsubsection{Description}
Sets the number of instructions to execute between checks for the pause character.

\subsection{Sim-CPU}
\subsubsection{Inputs}
A \datatype{String} containing the CPU name.
See the ``SET CPU'' (section \ref{cmd-set}) command for valid CPU names.
\subsubsection{Outputs}
None
\subsubsection{Description}
Selects the CPU to simulate.
This can only be done once and must be done before any command that attempts to access the CPU.

\subsection{Sim-init}
\subsubsection{Inputs}
None.
\subsubsection{Outputs}
None.
\subsubsection{Description}
Calls the initialization routine to the current simulator.
The action is simulator dependent, but typically clears the registers.
The memory may, or may not be cleared.

\subsection{Sim-load}
\subsubsection{Inputs}
A \datatype{String} containing a file name.
\subsubsection{Outputs}
None.
\subsubsection{Description}
Calls the load routine for the current simulator and passes the string containing the file name.
The effect is simulator dependent, but the 8080 simulator will attempt to load the file as an Intel hex format file and the 68000 simulator will attempt to load the file as a Motorola S-record file.
The PDP-11 simulator will load a binary file in absolute loader format.

\subsection{Sim-step}
\subsubsection{Inputs}
None.
\subsubsection{Outputs}
None.
\subsubsection{Description}
Executes one instruction or step of the simulator.

\subsection{Tape-Close}
The following inputs are expected for paper tape controllers:
\begin{itemize}
  \item A \datatype{String} representing the paper tape controller name.
  \item A \datatype{String} representing the output (``PUN'' for punch) or input (``RDR'' for reader).
\end{itemize}
The following inputs are expected for magnetic tape controllers:
\begin{itemize}
  \item A \datatype{String} representing the paper tape controller name.
  \item An \datatype{Integer} representing the tape drive to close.
\end{itemize}
\subsubsection{Outputs}
None
\subsubsection{Description}
Closes the file attached to the specified tape drive unit.

\subsection{Tape-Open}
The following inputs are expected for paper tape controllers:
\begin{itemize}
  \item A \datatype{String} representing the tape controller name.
  \item A \datatype{String} representing the output (``PUN'' for punch) or input (``RDR'' for reader).
  \item A \datatype{String} representing the file name to attach to the tape drive.
\end{itemize}
The following inputs are expected for magnetic tape controllers:
\begin{itemize}
  \item A \datatype{String} representing the tape controller name.
  \item An \datatype{Integer} representing the tape drive to open.
  \item A st\datatype{String}ring representing the file name to attach to the tape drive.
\end{itemize}
\subsubsection{Outputs}
None
\subsubsection{Description}
Opens a file and attaches it to the specified tape drive unit.

\subsection{Tape-protect}
The following inputs are expected for magnetic tape controllers:
\begin{itemize}
  \item A \datatype{String} representing the paper tape controller name.
  \item An \datatype{Integer} representing the tape drive to close.
  \item An \datatype{Integer} where a value of 0 sets the drive to read-write and any other value sets it to read-only.
\end{itemize}
\subsubsection{Outputs}
None
\subsubsection{Description}
Sets the read/write status of the specified tape drive unit.
This applies only to magnetic tape drives.

\subsection{Examples}
The Lisp environment is useful for writing automated tests and for writing utilities to enhance the CLI.
\subsubsection{Watch a Memory Location}
The following Lisp program will execute instructions until the specified memory location changes.
\lstset{language=[Tiny]Lisp}
\begin{lstlisting}
(defun watch (addr)
  (setq old-value (meml addr))
  (dowhile (= old-value (meml addr))
    (sim-step)))
\end{lstlisting}

%========================================================
\clearpage
\addcontentsline{toc}{chapter}{Bibleography}
\nocite{dec84}
\bibliographystyle{plain}
\bibliography{sim.bib}

\end{document}

